\Needspace{9\baselineskip}\begin{center}\textit{Two methods for distinguishing composite numbers from primes and finding the factors of the former.}\end{center}


\Needspace{6\baselineskip}\Article{329}
It is so well known that the problem of distinguishing primes from composite numbers and resolving the latter into their prime factors is among the most important and useful problems in all arithmetic, and has occupied the diligence and ingenuity of mathematicians both ancient and modern, that a lengthy discussion would be superfluous. Nevertheless, we must admit that all methods proposed so far are either restricted to very special cases or so laborious and lengthy that they exhaust even a practiced calculator's patience for numbers that do not exceed the limits of the tables constructed by deserving men, that is, numbers for which elaborate methods are unnecessary; for larger numbers they can generally hardly be used at all. Those tables, which are in everyone's hands and which we may hope will gradually be extended further, certainly suffice in most ordinary cases. Yet an experienced calculator quite often has occasion to gain great advantages from factoring large numbers, advantages that amply repay a moderate expenditure of time. Moreover, the dignity of the science seems to require that every means of solving so elegant and celebrated a problem be carefully developed. For these reasons we do not
\SourceLeaf{407}
doubt that the following two methods, whose effectiveness and brevity we can confirm from long experience, will be welcome to lovers of arithmetic. It is, moreover, inherent in the nature of the problem that \textit{every} method becomes progressively longer as the numbers to which it is applied grow larger. With the following methods, however, the difficulties increase very slowly, and numbers of seven, eight or even more digits have always been treated successfully, especially by the second method, with all the speed that can reasonably be expected for such large numbers, which by all methods known hitherto would require work intolerable even to an indefatigable calculator.

Before using the following methods, it is always very useful to try dividing the proposed number by some of the smallest primes, say \(2,\allowbreak 3,\allowbreak 5,\allowbreak 7\), etc., up to \(19\) or further. This is not only to avoid regretting the discovery, by subtle and elaborate methods, of a divisor that could have been found much more easily by division alone\HistoricalNote{*)}{All the more so because, generally speaking, scarcely one number in six is found not divisible by any of \(2,\allowbreak 3,\allowbreak 5\ldots19\).}, but also because, when none of these divisions succeeds, the second method makes very profitable use of the \textit{remainders} from them. Thus, for example, if \(314159265\) is to be resolved into its factors, division by \(3\) succeeds twice, followed by divisions by \(5\) and \(7\), giving \(314159265=9.5.7.997331\). It is therefore enough to examine more closely the number \(997331\), which is found not divisible by \(11,\allowbreak 13,\allowbreak 17,\allowbreak 19\). Similarly, given \(43429448\), we shall remove the factor \(8\) and apply the more elaborate methods to the quotient \(5428681\).

\Needspace{6\baselineskip}\Article{330}
The \textsc{first method} rests on the theorem that \textit{any positive or} \textit{negative number that is a quadratic residue of another number} \(M\) \textit{is also a residue of every} \textit{divisor of} \(M\)\textit{.} It is well known that, if \(M\) is divisible by no prime below \(\surd M\), then \(M\) is certainly prime. If, however, all primes below this bound dividing \(M\) are \(p,\allowbreak q\), etc., then \(M\) either consists of these \textit{alone}, or their powers, or can contain only \textit{one} other prime factor greater than \(\surd M\), which is found by dividing \(M\) by \(p,\allowbreak q\), etc., as often as possible. Denote the collection of all primes below \(\surd M\), excluding those by which division has already been tried unsuccessfully, by \(\Omega\). It is evident
\SourceLeaf{408}
that it is enough to find all prime divisors of \(M\) contained in \(\Omega\). Now, if we know from any source that a number \(r\), not a square, is a quadratic residue of \(M\), no prime for which \(r\) is \(NR.\) can be a divisor of \(M\). We may therefore remove from \(\Omega\) all such primes, which will generally make up about half of the whole. If another number that is not a square, \(r'\), is also known to be a residue of \(M\), we may again exclude from the primes remaining in \(\Omega\) after the first exclusion those for which \(r'\) is \(NR.\). These will again make up about half of them, provided the residues \(r\) and \(r'\) are independent, \textit{i.e.} unless one is necessarily, of itself, a residue of all numbers of which the other is a residue, as would happen if \(rr'\) were a square. If still other residues of \(M\) are known, \(r'',\allowbreak r'''\), etc., all independent of the others\HistoricalNote{*)}{If the product of any number of \(r,\allowbreak r',\allowbreak r''\), etc., is a square, each of them, for example \(r\), will be a residue of every prime, dividing none of these numbers, that is a residue of the remaining numbers \(r',\allowbreak r''\), etc. Thus, for any number of residues to be regarded as independent, no product of two, three, etc., may be a square.}, similar exclusions can be made with each. These reduce the collection of numbers in \(\Omega\) very rapidly, so that soon either all are removed, in which case \(M\) is certainly prime, or so few remain, evidently including all prime divisors of \(M\), if it has any, that division by them can be tried without difficulty. For a number not much above a million, six or seven exclusions will generally be more than enough; for a number of eight or nine digits, nine or ten. Two points now remain to be discussed: \textit{first}, how to find sufficiently many suitable residues of \(M\), and \textit{then}, how to perform the exclusion itself most conveniently. But we shall reverse the order of these questions, especially because the second will show which residues are most suitable for this purpose.
